\chapter{Перестановочный носитель спина, аромата и пространства}
\label{ch:v8-baryon-spin-flavor-permutation-carrier-gate}

Предыдущая глава условно использовала полностью симметричное
спин-ароматное состояние. Проверим, следует ли оно из цветового
эпсилон-синглета и принципа Паули без дополнительного пространственного
гамильтониана.

\section{Условие полной фермионной антисимметрии}

Цветовой эпсилон-вектор несёт знаковое представление группы перестановок:
\begin{equation}
 U_\pi^{\rm color}|\varepsilon\rangle
 =\operatorname{sgn}(\pi)|\varepsilon\rangle.
 \label{eq:v8-baryon-color-sign-representation}
\end{equation}
Поэтому полное трёхкварковое состояние антисимметрично тогда и только
тогда, когда произведение пространственной и спин-ароматной частей
содержит тривиальное представление $S_3$.

Для неприводимых типов
$\mathbf1$, $\mathbf1_{\rm sgn}$ и $\mathbf2$ их характеры на классах
$e$, $(12)$ и $(123)$ равны
$(1,1,1)$, $(1,-1,1)$ и $(2,0,-1)$. Кратность инварианта равна
\begin{equation}
 m_{\lambda\mu}
 =\frac16\left[
 \chi_\lambda(e)\chi_\mu(e)
 +3\chi_\lambda(12)\chi_\mu(12)
 +2\chi_\lambda(123)\chi_\mu(123)\right].
 \label{eq:v8-baryon-s3-invariant-multiplicity}
\end{equation}
Точный расчёт даёт
\begin{equation}
 (m_{\lambda\mu})_{\lambda,\mu\in
 \{\mathbf1,\mathbf1_{\rm sgn},\mathbf2\}}=I_3.
 \label{eq:v8-baryon-s3-matching-matrix}
\end{equation}
Следовательно, принцип Паули допускает три ветви:
\begin{equation}
 (\mathcal H_{\rm space},\mathcal H_{\rm sf})
 \sim(\mathbf1,\mathbf1),\quad
 (\mathbf1_{\rm sgn},\mathbf1_{\rm sgn}),\quad
 (\mathbf2,\mathbf2).
 \label{eq:v8-baryon-three-permutation-branches}
\end{equation}
Он требует совпадения перестановочных типов, но не выбирает тип
$\mathbf1$ по отдельности.

\section{Спин-ароматный сектор нуклона}

На носителе
$(\mathbb C^2_{\rm aroma}\otimes\mathbb C^2_{\rm spin})^{\otimes3}$
совместный сектор $I=S=1/2$ имеет размерность четыре и раскладывается как
\begin{equation}
 \mathcal H_{I=S=1/2}
 \cong\mathbf1\oplus\mathbf1_{\rm sgn}\oplus\mathbf2.
 \label{eq:v8-baryon-nucleon-s3-decomposition}
\end{equation}
Все три слагаемых совместимы с цветовым эпсилон-вектором, если
пространственная часть имеет тот же перестановочный тип.

Для зарядово-спинового оператора
$O=\sum_{r<s}Q_rQ_s\boldsymbol\Sigma_r\cdot\boldsymbol\Sigma_s$
получаем точную таблицу:
\begin{equation}
 \begin{array}{c|ccc}
 \text{тип }S_3&O_p&O_n&O_n-O_p\\ \hline
 \mathbf1&4/3&1&-1/3\\
 \mathbf1_{\rm sgn}&-4/3&-1/3&1\\
 \mathbf2&0&1/3&1/3
 \end{array}.
 \label{eq:v8-baryon-s3-magnetic-branch-table}
\end{equation}
Таким образом, все три фермионно-допустимые ветви дают разные величины, а
две из них меняют знак относительно полностью симметричной ветви.

\begin{theorem}[Неединственность перестановочного селектора]
Цветовая антисимметрия и принцип Паули не выбирают полностью симметричное
спин-ароматное состояние. Они выбирают только диагональное совпадение
перестановочных типов пространства и спина--аромата.
\end{theorem}

\section{Какое дополнительное условие было бы достаточно}

Если независимый пространственный родитель имеет единственное
перестановочно-симметричное основное состояние,
\begin{equation}
 \ker(H_{\rm space}-E_0)=\mathbf1,
 \label{eq:v8-baryon-symmetric-spatial-ground-condition}
\end{equation}
то принцип Паули действительно вынуждает
$\mathcal H_{\rm sf}\sim\mathbf1$ и восстанавливает условный множитель
$O_n-O_p=-1/3$. Но в текущей теории нет трёхкоординатного барионного
гамильтониана, его основного состояния или доказательства невырожденности.
Даже обозначение «$S$-волна» само по себе не заменяет это спектральное
условие.

\begin{nogocriterion}[Граница цветового эпсилон-проектора]
Эпсилон-проектор определяет знаковый цветовой тип, но не определяет
пространственный тип. Поэтому из него нельзя вывести ни полностью
симметричную спин-ароматную волну, ни магнитную разность $-1/3$ без
отдельного родителя пространственного основного состояния.
\end{nogocriterion}

\begin{protocol}
\begin{enumerate}
 \item характеры $S_3$: \textbf{проверены точно};
 \item матрица кратностей инварианта: \textbf{$I_3$};
 \item допустимые фермионные ветви: \textbf{три};
 \item размерность сектора $I=S=1/2$: \textbf{четыре};
 \item разложение: \textbf{$\mathbf1\oplus\mathbf1_{\rm sgn}\oplus\mathbf2$};
 \item магнитные значения трёх ветвей: \textbf{вычислены точно};
 \item селекция симметричной ветви принципом Паули: \textbf{запрещена};
 \item достаточное дополнительное условие: \textbf{единственное
 симметричное пространственное основное состояние};
 \item такое состояние в текущем проекте: \textbf{не выведено};
 \item следующий гейт: \textbf{симметрия пространственного основного
 состояния}.
\end{enumerate}
\end{protocol}

Машинный сертификат сохранён в
\path{s2t_v8_baryon_spin_flavor_permutation_carrier_gate_results.json}.